% TOP_PROBLEM: {"id": 3033, "title": "$C^r$ Stability Conjecture", "category_id": 8, "category": {"id": 8, "name": "analysis", "display_name": "Analysis", "description": "Limits, continuity, calculus, and function theory.", "slug": "analysis", "order_index": 8, "created_at": "2026-07-31T15:26:25.671Z"}, "status": "open", "problem_number": "OPG-725", "set_id": 12, "statusQualification": "Scope made explicit: closed manifolds and the standard Axiom A stability conjecture; hyperbolic does not mean Anosov on the whole manifold."}
% FIRST_POSTED: 2026-10-01
% ATTEMPT_STATUS: unresolved
% UnsolvedMath ID 3033; source catalogue code OPG-725
% !TeX program = lualatex
% Research attempt by GPT-6 Astra Ultra; no claim of a new complete solution.
\documentclass[11pt,a4paper]{article}
\usepackage[margin=25mm]{geometry}
\usepackage{fontspec}
\setmainfont{lmroman10-regular.otf}[BoldFont=lmroman10-bold.otf,
 ItalicFont=lmroman10-italic.otf,BoldItalicFont=lmroman10-bolditalic.otf]
\usepackage{amsmath,amssymb,mathtools}
\usepackage{unicode-math}
\setmathfont{latinmodern-math.otf}
\AtBeginDocument{\let\setminus\smallsetminus}
\usepackage{xurl,hyperref}
\hypersetup{colorlinks=true,urlcolor=blue}
\setcounter{secnumdepth}{0}
\setlength{\parindent}{0pt}
\setlength{\parskip}{5pt}
\setlength{\emergencystretch}{3em}
\begin{document}
\section*{$C^r$ Stability Conjecture}\label{3033}
{\small UnsolvedMath ID 3033 · OPG-725 · Analysis · OpenGarden}
\subsection{Definitions and mathematical statement}
Let $M$ be a smooth compact manifold without boundary, equipped with a
Riemannian norm, and let $r\in\{1,2,\ldots\}\cup\{\infty\}$.
A $C^r$ diffeomorphism $f:M\to M$ is $C^r$ structurally stable if some
neighbourhood $\mathcal U$ of $f$ in the $C^r$ topology has this property:
for every $g\in\mathcal U$ there exists a homeomorphism $h:M\to M$ satisfying
$h\circ f=g\circ h$. The conjugating map is not required to be differentiable.

The nonwandering set $\Omega(f)$ consists of points $x$ such that every
neighbourhood $U$ of $x$ satisfies $f^j(U)\cap U\ne\varnothing$ for some
integer $j\ge1$. It is hyperbolic if the tangent bundle over it has a continuous
$Df$-invariant splitting $E^s\oplus E^u$ and constants $C>0$, $0<\lambda<1$
such that, for $j\ge0$,
\[
 \|Df^j v\|\le C\lambda^j\|v\|\ (v\in E^s),\qquad
 \|Df^{-j}w\|\le C\lambda^j\|w\|\ (w\in E^u).
\]
The Axiom A condition additionally requires the periodic points, those with
$f^j(x)=x$ for some $j\ge1$, to be dense in $\Omega(f)$.

In its standard stability-conjecture formulation, the question is whether
every $C^r$ structurally stable diffeomorphism satisfies Axiom A, particularly
for $r>1$. This specifies the catalogue's abbreviated word ``hyperbolic'';
it does not demand a hyperbolic splitting over all of $M$ (the Anosov
condition). Both the topology of perturbations and the regularity of $f$
are fixed by $r$.
\subsection{Short English statement}
Does persistence up to topological conjugacy under small $C^r$ perturbations force hyperbolic nonwandering dynamics with dense periodic points?
\subsection{Research attempt}
\paragraph{Scope and process.}
This is a GPT-6 Astra Ultra research attempt dated 1 October 2026, using
primary-source browsing and elementary mathematical checks. The canonical
definition above is retained. Results below are restricted results or
reductions, not a claim of a new solution to the entire catalogue question.
No formal proof assistant or external mathematical referee checked this text.


\paragraph{Literature and logical scope.}
Ma\~n\'e's theorem [S3, R1] establishes that $C^1$ structural stability
implies Axiom A on closed manifolds. The $r>1$ formulation in [S2]
has a stronger topology of allowed perturbations; being stable against
that smaller collection does not imply $C^1$ stability. Our work below
constructs smooth instabilities and quantifies this topology obstruction.
It is not a reproof of Ma\~n\'e's theorem.

\paragraph{An explicit nonhyperbolic circle example.}
Work on $M=\mathbb R/\mathbb Z$. Fix $0<a<1/(2\pi)$ and take the lift
\[
 F(x)=x+a(1-\cos 2\pi x).
\]
It satisfies $F(x+1)=F(x)+1$ and
$F'(x)=1+2\pi a\sin2\pi x>0$, hence induces a smooth diffeomorphism
$f$. Its displacement lies in $[0,2a]\subset[0,1)$ and vanishes
exactly at integers, so $f$ has exactly one fixed point. Its multiplier
there is one, so that periodic point is not hyperbolic.

For $0<\delta<1-2a$ set $F_\delta(x)=F(x)+\delta$. The derivative
is unchanged and its displacement lies strictly between zero and one;
therefore the induced $f_\delta$ has no fixed point. These maps converge
to $f$ in every $C^r$ topology, including $C^\infty$. Topological
conjugacy maps the fixed-point set bijectively onto the fixed-point set,
because $h(f(x))=g(h(x))$. Thus $f$ and $f_\delta$ cannot be conjugate.
This explicitly proves instability in all regularities for this parabolic
example, with no appeal to a perturbation lemma.

\paragraph{A neutral recurrent example without fixed points.}
For a rigid circle rotation $R_\alpha(x)=x+\alpha$, the multiplier
is one everywhere. If $\alpha$ is rational, every point is periodic;
if irrational, no point is periodic, since $q\alpha\notin\mathbb Z$
for every positive $q$. Rational and irrational parameters approximate
one another arbitrarily closely, and changing the parameter gives
convergence in every $C^r$ norm on lifts. Existence of a periodic point
is a conjugacy invariant, so every rigid rotation is structurally
unstable in every specified regularity. This checks a different
mechanism from the isolated parabolic fixed point above.

\paragraph{Why a $C^1$ perturbation proof does not automatically lift.}
In a local coordinate interval let $\phi$ be a nonzero smooth bump
supported in $(-1,1)$, and consider a displacement
$b_\varepsilon(x)=\varepsilon^2\phi(x/\varepsilon)$.
Its $j$th derivative has supremum
$\varepsilon^{2-j}\|\phi^{(j)}\|_\infty$.
Thus $b_\varepsilon\to0$ in $C^1$, but its second derivative stays
of fixed nonzero size, and higher derivatives may diverge. Extending
by zero gives smooth local perturbations; for small $\varepsilon$
the derivative change is small enough that a diffeomorphism remains
a diffeomorphism. These are legitimate small $C^1$ perturbations
which are not small $C^2$ perturbations. A proof that relies on their
availability has not met the $r=2$ hypothesis.

\paragraph{Verification and remaining gap.}
The fixed-point count used above is on the quotient circle, so we checked
that the displacement cannot equal a nonzero integer. Merely checking
that it is nonzero as a real number would not have sufficed. The
rotation argument uses periodic points of every period, not a change
of a chosen numerical rotation parameter alone, which conjugacy could
otherwise obscure. The bump calculation makes no claim that every
nonhyperbolic orbit can be handled by this particular bump. In higher
dimensions the missing step is a perturbation theorem, small in the
fixed $C^r$ topology, which turns general failure of Axiom A into a
topological conjugacy obstruction. Neither circle construction supplies
that theorem or proves density of periodic points in the general
nonwandering set.

\paragraph{Final status.}
\textbf{UNRESOLVED} for the general higher-regularity target. The
$C^1$ case is attributed to existing literature; two explicit smooth
families and the precise norm obstruction are independently verified here.

\subsection{Sources}
\begin{itemize}
\item[S1] ULAM.AI and the UnsolvedMath Contributors, supplied problems.json, record 3033 (OPG-725), OpenGarden collection. Catalogue identity and source transcription; CC BY 4.0.
\url{https://huggingface.co/datasets/ulamai/UnsolvedMath}.
\item[S2] Open Problem Garden, $C^r$ Stability Conjecture. Original problem and discussion; the mathematical formulation below is expanded from the supplied source record.
\url{http://www.openproblemgarden.org/op/c_r_stability_conjecture}.
\item[S3] R. Mañé, A proof of the C1 stability conjecture, Publications Mathématiques de l’IHÉS 66 (1987), 161–210, introduction and definitions.
\url{https://www.numdam.org/article/PMIHES_1987__66__161_0.pdf}.

\item[R1] Ricardo Ma\~n\'e, \emph{A proof of the $C^1$ stability
conjecture}, Publications Math\'ematiques de l'IH\'ES 66 (1987),
161--210. Primary publisher archive, checked 1 October 2026.
\url{https://numdam.org/articles/10.1007/BF02698931/}.

\end{itemize}
\end{document}
